\ifdefined\gkver\else\def\gkver{student}\fi
\documentclass[\gkver]{gaokaozhenti}
\begin{document}

\chapter{2009年天津理综}
%% paperType: 理综物理部分

\begin{enumerate}
\item
%% number: 1
%% paperName: 2009年天津理综第 1 题 6 分
%% typeId: 1
%% type: 单选题
%% chapter: 力物体的平衡
%% point: 多力平衡
%% method: 无
%% score: 6
%% degree: 650
%% duplicateId: 0
%% body:
物块静止在固定的斜面上，分别按图示的方向对物块施加大小相等的力 $F$，A 中 $F$ 垂直于斜面向上，B 中 $F$ 垂直于斜面向下，C 中 $F$ 竖直向上，D 中 $F$ 竖直向下，施力后物块仍然静止，则物块所受的静摩擦力增大的是\xzanswer{D}

\fourchoices[answer=D, ispicture=true, h=1.9cm]
{figs/01a.png}
{figs/01b.png}
{figs/01c.png}
{figs/01d.png}

%% answer: D
%% memo:
\memoanswer{
四个图中都是静摩擦。A 图中 $F_{fA}=G\sin\theta$；B 图中 $F_{fB}=G\sin\theta$；C 图中 $F_{fC}=(G-F)\sin\theta$；D 图中 $F_{fC}=(G+F)\sin\theta$。

\onepicture[w=13cm]{figs/01_an.png}
}

\item
%% number: 2
%% paperName: 2009年天津理综第 2 题 6 分
%% typeId: 1
%% type: 单选题
%% chapter: 光学
%% point: 光的电磁说
%% method: 无
%% score: 6
%% degree: 650
%% duplicateId: 0
%% body:
下列关于电磁波的说法正确的是\xzanswer{B}

\fourchoices[answer=B]
{电磁波必须依赖介质传播}
{电磁波可以发生衍射现象}
{电磁波不会发生偏振现象}
{电磁波无法携带信息传播}

%% answer: B
%% memo:
\memoanswer{
电磁波在真空中也能传播，A 错；衍射是一切波所特有的现象，B 对；电磁波是横波，横波能发生偏振现象，C 错；所有波都能传递信息，D 错。
}

\item
%% number: 3
%% paperName: 2009年天津理综第 3 题 6 分
%% typeId: 1
%% type: 单选题
%% chapter: 电路
%% point: 电路中的图象
%% method: 无
%% score: 6
%% degree: 650
%% duplicateId: 0
%% body:
为探究小灯泡 L 的伏安特性，连好图示的电路后闭合开关，通过移动变阻器的滑片，使小灯泡中的电流由零开始逐渐增大，直到小灯泡正常发光。由电流表和电压表得到的多组读数描绘出的 $U-I$ 图象应是\xzanswer{C}

\onepicture[w=4.2cm]{figs/03.png}

\fourchoices[answer=C, ispicture=true, w=2.6cm]
{figs/03a.png}
{figs/03b.png}
{figs/03c.png}
{figs/03d.png}

%% answer: C
%% memo:
\memoanswer{
灯丝电阻随电压的增大而增大，在图像上某点到原点连线的斜率应越来越大。C 正确。
}

\item
%% number: 4
%% paperName: 2009年天津理综第 4 题 6 分
%% typeId: 1
%% type: 单选题
%% chapter: 电磁感应
%% point: 电磁感应能量分析
%% method: 无
%% score: 6
%% degree: 450
%% duplicateId: 0
%% body:
如图所示，竖直放置的两根平行金属导轨之间接有定值电阻 $R$，质量不能忽略的金属棒与两导轨始终保持垂直并良好接触且无摩擦，棒与导轨的电阻均不计，整个装置放在匀强磁场中，磁场方向与导轨平面垂直，棒在竖直向上的恒力 $F$ 作用下加速上升的一段时间内，力 $F$ 做的功与安培力做的功的代数和等于\xzanswer{A}

\onepicture[w=2.6cm]{figs/04.png}

\fourchoices[answer=A]
{棒的机械能增加量}
{棒的动能增加量}
{棒的重力势能增加量}
{电阻 $R$ 上放出的热量}

%% answer: A
%% memo:
\memoanswer{
棒受重力 $G$、拉力 $F$ 和安培力 $F_{A}$ 的作用。由动能定理：$W_{F}+W_{G}+W_{\text{安}}=\Delta E_{K}$ 得 $W_{F}+W_{\text{安}}=\Delta E_{K}+mgh$，即力 $F$ 做的功与安培力做功的代数和等于机械能的增加量。选 A。
}

\item
%% number: 5
%% paperName: 2009年天津理综第 5 题 6 分
%% typeId: 1
%% type: 单选题
%% chapter: 电场
%% point: 带电粒子的直线运动
%% method: 无
%% score: 6
%% degree: 450
%% duplicateId: 0
%% body:
如图所示，带等量异号电荷的两平行金属板在真空中水平放置，M、N 为板间同一电场线上的两点，一带电粒子（不计重力）以速度 $v_{M}$ 经过 M 点在电场线上向下运动，且未与下板接触，一段时间后，粒子以速度 $v_{N}$ 折回 N 点。则\xzanswer{B}

\onepicture[w=4.5cm]{figs/05.png}

\fourchoices[answer=B]
{粒子受电场力的方向一定由 M 指向 N}
{粒子在 M 点的速度一定比在 N 点的大}
{粒子在 M 点的电势能一定比在 N 点的大}
{电场中 M 点的电势一定高于 N 点的电势}

%% answer: B
%% memo:
\memoanswer{
由于带电粒子未与下板接触，可知粒子向下做的是减速运动，故电场力向上，A 错；粒子由 M 到 N 电场力做负功电势能增加，动能减少，速度增加，故 B 对 C 错；由于粒子和两极板所带电荷的电性未知，故不能判断 M、N 点电势的高低，C 错。
}

\item
%% number: 6
%% paperName: 2009年天津理综第 6 题 6 分
%% typeId: 2
%% type: 多选题
%% chapter: 原子和原子核
%% point: 人工核反应
%% method: 无
%% score: 6
%% degree: 650
%% duplicateId: 0
%% body:
下列说法正确的是\xzanswer{BD}

\fourchoices[answer=BD]
{$\ce{^{15}_{7}N + ^{1}_{1}H -> ^{12}_{6}C + ^{4}_{2}He}$ 是 $\alpha$ 衰变方程}
{$\ce{^{1}_{1}H + ^{2}_{1}H -> ^{3}_{2}He + \gamma}$ 是核聚变反应方程}
{$\ce{^{238}_{92}U -> ^{234}_{90}Th + ^{4}_{2}He}$ 是核裂变反应方程}
{$\ce{^{4}_{2}He + ^{27}_{13}Al -> ^{30}_{15}P + ^{1}_{0}n}$ 是原子核的人工转变方程}

%% answer: BD
%% memo:
\memoanswer{
A 选项中 $\ce{^{15}_{7}N}$ 在质子的轰击下发生的核反应，属于人工转变，A 错；C 选项是 $\alpha$ 衰变，不是裂变，C 错。
}

\item
%% number: 7
%% paperName: 2009年天津理综第 7 题 6 分
%% typeId: 2
%% type: 多选题
%% chapter: 光学
%% point: 光的折射
%% method: 无
%% score: 6
%% degree: 450
%% duplicateId: 0
%% body:
已知某玻璃对蓝光的折射率比对红光的折射率大，则两种光\xzanswer{C}

\fourchoices[answer=C]
{在该玻璃中传播时，蓝光的速度较大}
{以相同的入射角从空气斜射入该玻璃中，蓝光折射角较大}
{从该玻璃中射入空气发生反射时，红光临界角较大}
{用同一装置进行双缝干涉实验，蓝光的相邻条纹间距较大}

%% answer: C
%% memo:
\memoanswer{
由 $v=\frac{c}{n}$ 可知，蓝光在玻璃中的折射率大，蓝光的速度较小，A 错；以相同的入射角从空气中斜射入玻璃中，蓝光的折射率大，向法线靠拢偏折得多，折射角应较小，B 错。从玻璃射入空气发生全反射时的临界角由公式 $\sin C=\frac{1}{n}$ 可知，红光的折射率小，临界角大，C 正确；用同一装置进行双缝干涉实验，由公式 $\Delta x=\frac{L}{d}\lambda$ 可知蓝光的波长短，相邻条纹间距小，D 错。
}

\item
%% number: 8
%% paperName: 2009年天津理综第 8 题 6 分
%% typeId: 2
%% type: 多选题
%% chapter: 机械振动
%% point: 简谐振动
%% method: 无
%% score: 6
%% degree: 450
%% duplicateId: 0
%% body:
某质点做简谐运动，其位移随时间变化的关系式为 $x=A\sin\frac{\pi}{4}t$，则质点\xzanswer{AD}

\fourchoices[answer=AD]
{第 $1\Us$ 末与第 $3\Us$ 末的位移相同}
{第 $1\Us$ 末与第 $3\Us$ 末的速度相同}
{$3\Us$ 末至 $5\Us$ 末的位移方向都相同}
{$3\Us$ 末至 $5\Us$ 末的速度方向都相同}

%% answer: AD
%% memo:
\memoanswer{
由关系式可知 $\omega=\frac{\pi}{4}\Urads$，$T=\frac{2\pi}{\omega}=8\Us$，将 $t=1\Us$ 和 $t=3\Us$ 代入关系式中求得两时刻位移相同，A 对；画出对应的位移-时间图像，由图像可以看出，第 $1\Us$ 末和第 $3\Us$ 末的速度方向不同，B 错；仍由图像可知，$3\Us$ 末至 $5\Us$ 末的位移大小相同，方向相反，而速度是大小相同，方向也相同。故 C 错、D 对。

\onepicture[w=7cm]{figs/08_an.png}
}

\item
%% number: 9
%% paperName: 2009年天津理综第 9 题 8 分
%% typeId: 4
%% type: 实验题
%% chapter: 直线运动
%% point: 自由落体运动
%% method: 无
%% score: 8
%% degree: 450
%% duplicateId: 0
%% body:
\begin{enumerate}
	\item
	如图所示，单匝矩形闭合导线框 abcd 全部处于磁感应强度为 $B$ 的水平匀强磁场中，线框面积为 $S$，电阻为 $R$。线框绕与 cd 边重合的竖直固定转轴以角速度 $\omega$ 匀速转动，线框中感应电流的有效值 $I=$ \tkanswer[1.6cm]{}。线框从中性面开始转过 $\frac{\pi}{2}$ 的过程中，通过导线横截面的电荷量 $q=$ \tkanswer[1.6cm]{}。

	\onepicture[align=right, w=3.0cm]{figs/09a.png}

	\item
	图示为简单欧姆表原理示意图，其中电流表的偏电流 $I_{R}=300\UuA$，内阻 $R_{g}=100\UO$，可变电阻 $R$ 的最大阻值为 $10\UkO$，电池的电动势 $E=1.5\UV$，内阻 $r=0.5\UO$，图中与接线柱 A 相连的表笔颜色应是 \tkanswer[1.6cm]{} 色，接正确使用方法测量电阻 $R_{x}$ 的阻值时，指针指在刻度盘的正中央，则 $R_{x}=$ \tkanswer[1.2cm]{} $\UkO$。若该欧姆表使用一段时间后，电池电动势变小，内阻变大，但此表仍能调零，按正确使用方法再测上述 $R_{x}$ 其测量结果与原结果相比较 \tkanswer[1.6cm]{}（填“变大”、“变小”或“不变”）。

	\onepicture[align=right, w=3.2cm]{figs/09b.png}

	\item
	如图所示，将打点计时器固定在铁架台上，使重物带动纸带从静止开始自由下落，利用此装置可以测定重力和速度。

	\onepicture[align=right, w=4.2cm]{figs/09c.png}

	\begin{steps}
		\item
		所需器材有打点计时器（带导线）、纸带、复写纸、带铁夹的铁架台和带夹子的重物，此外还需 \tkanswer[1.6cm]{}（填字母代号）中的器材。

		（A）直流电源、天平及砝码

		（B）直流电源、毫米刻度尺

		（C）交流电源、天平及砝码

		（D）交流电源、毫米刻度尺
		\item
		通过作图象的方法可以剔除偶然误差较大的数据，提高实验的准确程度。为使图线的斜率等于重力加速度，除作 $v-t$ 图象外，还可作 \tkanswer[2.0cm]{} 图象，其纵轴表示的是 \tkanswer[2.4cm]{}，横轴表示的是 \tkanswer[2.4cm]{}。
	\end{steps}
\end{enumerate}
%% answer: 
\jdanswer{
\begin{enumerate}
	\item
	$\frac{\sqrt{2}BS\omega}{2R}$，$\frac{BS}{R}$

	\item
	红，5，变大

	\item
	①D；②$\frac{v^{2}}{2}-h$，速度平方的二分之一，重物下落的高度。
\end{enumerate}
}
%% memo:
\memoanswer{
（1）本题考查交变电流的产生和最大值、有效值、平均值的关系及交变电流中有关电荷量的计算等知识。电动势的最大值 $E_{m}=BS\omega$，电动势的有效值 $E=\frac{E_{m}}{\sqrt{2}}$，电流的有效值 $I=\frac{E}{R}=\frac{\sqrt{2}BS\omega}{2R}$；$q=\overline{I}\Delta t=\frac{\overline{E}}{R}\Delta t=\frac{\Delta\Phi}{R\Delta t}\Delta t=\frac{\Delta\Phi}{R}=\frac{BS}{R}$。

（2）本题考查欧姆表的结构、测量原理和相关计算，还涉及测量误差的分析。欧姆表是电流表改装的，必须满足电流的方向“+”进“-”出，即回路中电流从标有“+”标志的红表笔进去，所以与 A 相连的表笔颜色是红色；当两表笔短接（即 $R_{x}=0$）时，电流表应调至满偏电流 $I_{g}$，设此时欧姆表的内阻为 $R_{\text{内}}$，此时有关系 $I_{g}=\frac{E}{R_{\text{内}}}$ 得 $R_{\text{内}}=\frac{E}{I_{g}}=5\UkO$；当指针指在刻度盘的正中央时 $I=I_{g}/2$，有 $\frac{I_{g}}{2}=\frac{E}{R_{\text{内}}+R_{x}}$，代入数据可得 $R_{x}=R_{\text{内}}=5\UkO$；当电池电动势变小、内阻变大时，欧姆表得重新调零，由于满偏电流 $I_{g}$ 不变，由公式 $I_{g}=\frac{E}{R_{\text{内}}}$，欧姆表内阻 $R_{\text{内}}$ 得调小，待测电阻的测量值是通过电流表的示数体现出来的，由 $I=\frac{E}{R_{\text{内}}+R_{x}}=\frac{I_{g}R_{\text{内}}}{R_{\text{内}}+R_{x}}=\frac{I_{g}}{1+\frac{R_{x}}{R_{\text{内}}}}$，可知当 $R_{\text{内}}$ 变小时，$I$ 变小，指针跟原来的位置相比偏左了，欧姆表的示数变大了。

（3）本题考查用打点计时器测重力加速度。涉及器材的选取和用图像处理数据的方法。①打点计时器需接交流电源。重力加速度与物体的质量无关，所以不要天平和砝码。计算速度需要测相邻计数的距离，需要刻度尺，选 D。②由公式 $v^{2}=2gh$，如绘出 $\frac{v^{2}}{2}-h$ 图像，其斜率也等于重力加速度。
}

\item
%% number: 10
%% paperName: 2009年天津理综第 10 题 16 分
%% typeId: 6
%% type: 计算题
%% chapter: 运动和力
%% point: 动量和能量
%% method: 无
%% score: 16
%% degree: 450
%% duplicateId: 0
%% body:
如图所示，质量 $m_{1}=0.3\Ukg$ 的小车静止在光滑的水平面上，车长 $L=15\Um$，现有质量 $m_{2}=0.2\Ukg$ 可视为质点的物块，以水平向右的速度 $v_{0}=2\Ums$ 从左端滑上小车，最后在车面上某处与小车保持相对静止。物块与车面间的动摩擦因数 $\mu=0.5$，取 $g=10\Umsq$。求：

\begin{enumerate}
	\item
	物块在车面上滑行的时间 $t$；
	\item
	要使物块不从小车右端滑出，物块滑上小车左端的速度 $v_{0}'$ 不超过多少。
\end{enumerate}
\onepicture[align=right, w=5.0cm]{figs/10.png}
%% answer: 
\jdanswer{
\begin{enumerate}
	\item
	$0.24\Us$

	\item
	$5\Ums$
\end{enumerate}
}
%% memo:
\memoanswer{
本题考查摩擦拖动类的动量和能量问题。涉及动量守恒定律、动量定理和功能关系这些物理规律的运用。

（1）设物块与小车的共同速度为 $v$，以水平向右为正方向，根据动量守恒定律有 $m_{2}v_{0}=(m_{1}+m_{2})v$ ①。设物块与车面间的滑动摩擦力为 $F$，对物块应用动量定理有 $-Ft=m_{2}v-m_{2}v_{0}$ ②，其中 $F=\mu m_{2}g$ ③。解得 $t=\frac{m_{1}v_{0}}{\mu(m_{1}+m_{2})g}$，代入数据得 $t=0.24\Us$ ④。

（2）要使物块恰好不从车厢滑出，须物块到车面右端时与小车有共同的速度 $v'$，则 $m_{2}v_{0}'=(m_{1}+m_{2})v'$ ⑤。由功能关系有 $\frac{1}{2}m_{2}v_{0}'^{2}=\frac{1}{2}(m_{1}+m_{2})v'^{2}+\mu m_{2}gL$ ⑥。代入数据解得 $v_{0}'=5\Ums$。故要使物块不从小车右端滑出，物块滑上小车的速度 $v_{0}'$ 不能超过 $5\Ums$。
}

\item
%% number: 11
%% paperName: 2009年天津理综第 11 题 18 分
%% typeId: 6
%% type: 计算题
%% chapter: 磁场
%% point: 带电粒子在复合场中的运动
%% method: 无
%% score: 18
%% degree: 450
%% duplicateId: 0
%% body:
如图所示，直角坐标系 $xOy$ 位于竖直平面内，在水平的 $x$ 轴下方存在匀强磁场和匀强电场，磁场的磁感应强度为 $B$，方向垂直 $xOy$ 平面向里，电场线平行于 $y$ 轴。一质量为 $m$、电荷量为 $q$ 的带正电的小球，从 $y$ 轴上的 A 点水平向右抛出，经 $x$ 轴上的 M 点进入电场和磁场，恰能做匀速圆周运动，从 $x$ 轴上的 N 点第一次离开电场和磁场，MN 之间的距离为 $L$，小球过 M 点时的速度方向与 $x$ 轴的方向夹角为 $\theta$。不计空气阻力，重力加速度为 $g$，求：

\begin{enumerate}
	\item
	电场强度 $E$ 的大小和方向；
	\item
	小球从 A 点抛出时初速度 $v_{0}$ 的大小；
	\item
	A 点到 $x$ 轴的高度 $h$。
\end{enumerate}
\onepicture[align=right, w=3.6cm]{figs/11.png}
%% answer: 
\jdanswer{
\begin{enumerate}
	\item
	$\frac{mg}{q}$，方向竖直向上

	\item
	$\frac{qBL}{2m}\cot\theta$

	\item
	$\frac{q^{2}B^{2}L^{2}}{8m^{2}g}$
\end{enumerate}
}
%% memo:
\memoanswer{
本题考查平抛运动和带电小球在复合场中的运动。

（1）小球在电场、磁场中恰能做匀速圆周运动，说明电场力和重力平衡（恒力不能充当圆周运动的向心力），有 $qE=mg$ ①，$E=\frac{mg}{q}$ ②。重力的方向竖直向下，电场力方向只能向上，由于小球带正电，所以电场强度方向竖直向上。

（2）小球做匀速圆周运动，$O'$ 为圆心，MN 为弦长，$\angle MO'P=\theta$，如图所示。设半径为 $r$，由几何关系知 $\frac{L}{2r}=\sin\theta$ ③。小球做匀速圆周运动的向心力由洛伦兹力提供，设小球做圆周运动的速率为 $v$，有 $qvB=\frac{mv^{2}}{r}$ ④。由速度的合成与分解知 $\frac{v_{0}}{v}=\cos\theta$ ⑤。由③④⑤式得 $v_{0}=\frac{qBL}{2m}\cot\theta$ ⑥。

（3）设小球到 M 点时的竖直分速度为 $v_{y}$，它与水平分速度的关系为 $v_{y}=v_{0}\tan\theta$ ⑦。由匀变速直线运动规律 $v_{y}^{2}=2gh$ ⑧。由⑥⑦⑧式得 $h=\frac{q^{2}B^{2}L^{2}}{8m^{2}g}$ ⑨。

\onepicture[w=5.5cm]{figs/11_an.png}
}

\item
%% number: 12
%% paperName: 2009年天津理综第 12 题 20 分
%% typeId: 6
%% type: 计算题
%% chapter: 曲线运动
%% point: 天体运动
%% method: 无
%% score: 20
%% degree: 250
%% duplicateId: 0
%% body:
2008 年 12 月，天文学家们通过观测的数据确认了银河系中央的黑洞“人马座 A*”的质量与太阳质量的倍数关系。研究发现，有一星体 S2 绕人马座 A* 做椭圆运动，其轨道半长轴为 $9.50\times10^{2}$ 天文单位（地球公转轨道的半径为一个天文单位），人马座 A* 就处在该椭圆的一个焦点上。观测得到 S2 星的运行周期为 15.2 年。

\begin{enumerate}
	\item
	若将 S2 星的运行轨道视为半径 $r=9.50\times10^{2}$ 天文单位的圆轨道，试估算人马座 A* 的质量 $M_{A}$ 是太阳质量 $M_{s}$ 的多少倍（结果保留一位有效数字）；
	\item
	黑洞的第二宇宙速度极大，处于黑洞表面的粒子即使以光速运动，其具有的动能也不足以克服黑洞对它的引力束缚。由于引力的作用，黑洞表面处质量为 $m$ 的粒子具有势能为 $E_{p}=-G\frac{Mm}{R}$（设粒子在离黑洞无限远处的势能为零），式中 $M$、$R$ 分别表示黑洞的质量和半径。已知引力常量 $G=6.7\times10^{-11}\UNmqkgq$，光速 $c=3.0\times10^{8}\Ums$，太阳质量 $M_{s}=2.0\times10^{30}\Ukg$，太阳半径 $R_{s}=7.0\times10^{8}\Um$，不考虑相对论效应，利用上问结果，在经典力学范围内求人马座 A* 的半径 $R_{A}$ 与太阳半径 $R_{s}$ 之比应小于多少（结果按四舍五入保留整数）。
\end{enumerate}
%% answer: 
\jdanswer{
\begin{enumerate}
	\item
	$4\times10^{6}$

	\item
	$<17$
\end{enumerate}
}
%% memo:
\memoanswer{
本题考查天体运动的知识。其中第 2 小题为信息题，如“黑洞”“引力势能”等陌生的知识都在题目中给出，考查学生提取信息、处理信息的能力，体现了能力立意。

（1）S2 星绕人马座 A* 做圆周运动的向心力由人马座 A* 对 S2 星的万有引力提供，设 S2 星的质量为 $m_{S2}$，角速度为 $\omega$，周期为 $T$，则 $G\frac{M_{A}m_{S2}}{r^{2}}=m_{S2}\omega^{2}r$ ①，$\omega=\frac{2\pi}{T}$ ②。设地球质量为 $m_{E}$，公转轨道半径为 $r_{E}$，周期为 $T_{E}$，则 $G\frac{M_{S}m_{E}}{r_{E}^{2}}=m_{E}\omega^{2}r_{E}$ ③。综合上述三式得 $\frac{M_{A}}{M_{S}}=(\frac{r}{r_{E}})^{3}(\frac{T_{E}}{T})^{2}$。式中 $T_{E}=1$ 年 ④，$r_{E}=1$ 天文单位 ⑤。代入数据可得 $\frac{M_{A}}{M_{S}}=4\times10^{6}$ ⑥。

（2）引力对粒子作用不到的地方即为无限远，此时粒子的势能为零。“处于黑洞表面的粒子即使以光速运动，其具有的动能也不足以克服黑洞对它的引力束缚”，说明了黑洞表面处以光速运动的粒子在远离黑洞的过程中克服引力做功，粒子在到达无限远之前，其动能便减小为零，此时势能仍为负值，则其能量总和小于零，则有 $\frac{1}{2}mc^{2}-G\frac{Mm}{R}<0$ ⑦。依题意可知 $R=R_{A}$，$M=M_{A}$，可得 $R_{A}<\frac{2GM_{A}}{c^{2}}$ ⑧。代入数据得 $R_{A}<1.2\times10^{10}\Um$ ⑨，$\frac{R_{A}}{R_{S}}<17$ ⑩。
}

\end{enumerate}

\end{document}
